\documentclass[11pt]{article}

\usepackage[margin=0.85in]{geometry}
\usepackage{graphicx}
\usepackage{booktabs}
\usepackage{microtype}
\usepackage[hidelinks]{hyperref}
\usepackage{titlesec}
\usepackage{xcolor}
\usepackage{caption}

\titleformat{\section}[runin]{\bfseries\large}{\thesection.}{0.5em}{}[.]
\titleformat{\subsection}[runin]{\bfseries}{\thesubsection.}{0.5em}{}[.]
\titlespacing*{\section}{0pt}{0.6em}{0.5em}
\titlespacing*{\subsection}{0pt}{0.4em}{0.4em}
\setlength{\parskip}{0.25em}
\setlength{\parindent}{0pt}
\captionsetup{font=small,labelfont=bf,skip=2pt}

\title{\vspace{-2em}\textbf{Hierarchical Coordination for Real-Time
Multi-Modal Microtransit Dispatch:\\
Preliminary Evidence for NSF TTP-T Thrust 2}\vspace{-0.5em}}
\author{}
\date{}

\begin{document}
\maketitle
\vspace{-3em}

\section{Motivation and Research Question}
Thrust~2 (\emph{Robust Real-Time Operations for Workforce Commute}) targets
the dispatching layer of a multi-modal rural transit system. The dispatcher
must place hundreds to thousands of on-demand vehicles against time-windowed,
shift-synchronized demand under stochastic travel times. RQ~5 explicitly
proposes \textit{``hierarchical decomposition together with
reinforcement-learning techniques''} to scale to realistic networks. This
white paper reports preliminary evidence that hierarchical coordination is
viable as the operational backbone of Thrust~2, identifies which sub-problem
admits learning-based gains at v1 scale, and frames how the result composes
with Thrust~2's uncertainty-aware traffic models (Barrier~1) and the
digital-twin evaluation environment (Barrier~3).

\section{Approach}
We instantiate RQ~5's hierarchical control framework as a two-level
\textbf{HLP/LLP} decomposition. The \textbf{High-Level Planner} (HLP)
periodically partitions the fleet across hub-catchment regions on a forecast
window. The \textbf{Low-Level Planner} (LLP) performs per-region greedy
insertion of incoming requests, subject to capacity, time-window, and
return-to-depot feasibility. Decomposition is over the hub set rather than
$k$-means clusters of historical demand, because workforce commute is
destination-anchored.

The substrate is a real Memphis OSM drive network (5 hubs, up to 65 zones,
$\le$1{,}019 requests/day) with a perfect-information demand oracle as the
deterministic upper bound. Dispatch wall-time and on-time arrival are
measured against a flat \textbf{monolithic} baseline. Both planners are
configurable as \textit{heuristic} or \textit{learned}, so all four
combinations are sweepable from a single Hydra config.

\section{Preliminary Results}
\textbf{(i) Hierarchical decomposition alone delivers a clean speedup.} A
hub-catchment HLP plus heuristic greedy LLP gives a stable
$3.0\text{--}4.0\times$ per-decision wall-time speedup over the monolithic
baseline across fleet$\times$intensity grids on both 30- and 65-zone
substrates, with quality at parity at saturation and \emph{better} at
constrained corners (anti-vampire effect: regions stop one greedy decision
from draining a far-away vehicle; +12.8\,pp at fleet=60, intensity=3, K=5).

\textbf{(ii) A learned LLP compounds with hierarchical decomposition.}
Trained as a two-headed MLP (feasibility logit + log-cost regression) on
390k (vehicle, request) candidate evaluations, the model achieves
$98.9\%$ feasibility accuracy and $0.52$\,min cost MAE. The naive
\textit{rank-mode} plug-in (top-$K$ by predicted cost) under-performs at
saturation due to cost-ranking error; using only the model's reliable
feasibility head as a \textit{filter} (drop candidates with
$P(\text{feas.})<0.018$, exhaustive $(p,q)$ on survivors) preserves quality.
\textit{Stacked-mode} (filter then top-$K$ by cost on survivors) recovers
additional speedup. Compounded factors versus \textit{mono+heur} baseline
are reported in Table~\ref{tab:speedup}; Fig.~\ref{fig:wall} shows the full
speed--quality picture.

\textbf{(iii) A learned HLP yields insufficient signal at v1 scale.}
Spot-checking the heuristic HLP allocation against random
$\pm$5-vehicle perturbations (re-running the day with the
override at one rebalance tick) found only $\sim$27\% of ticks beatable
by $>$0.5\,pp, with peak gain $+3.14$\,pp at a single constrained-corner
tick and $\approx$0\,pp at light or moderate loads. With 5 hubs and a
perfect-information demand oracle, the proportional-to-demand heuristic
is near-optimal: there is little room to learn. Learned HLP is parked as a
v2 axis behind the multi-modal extension, where mode-mix and stochastic
travel will materially expand the action space; the data-collection
infrastructure is in place for cheap re-entry.

\begin{table}[ht]
\small
\centering
\caption{Speedup factor (mean per-decision wall-time, vs.\ \textit{mono+heur}
baseline) and on-time arrival delta (percentage points) at three
representative cells on the 30-zone Memphis OSM substrate, mean over
3 seeds. Filter mode preserves quality across the grid; stacked mode is
the upper bound at the cost of partial saturation regret.}
\label{tab:speedup}
\begin{tabular}{lccc cccc}
\toprule
& \multicolumn{3}{c}{\textbf{Speedup factor (vs.\ baseline)}}
& & \multicolumn{3}{c}{\textbf{$\Delta$on-time (pp)}} \\
\cmidrule(lr){2-4} \cmidrule(lr){6-8}
\textbf{Operating cell} & hier+heur & hier+filter & hier+stacked
& & hier+filter & hier+stacked \\
\midrule
fleet=120, int=3 (moderate)
& 3.59$\times$ & 8.23$\times$ & 9.92$\times$ & & $-3.3$ & $-7.1$ \\
fleet=120, int=5 (saturated)
& 3.70$\times$ & 8.06$\times$ & 8.98$\times$ & & $+3.5$ & $-12.1$ \\
fleet=240, int=5 (heavy load) %
& 3.81$\times$ & \textbf{11.67}$\times$ & \textbf{17.90}$\times$ & & $0.0$ & $-1.4$ \\
\bottomrule
\end{tabular}
\end{table}

\begin{figure}[ht]
\centering
\includegraphics[width=0.95\textwidth]{../learned_llp/wall_vs_fleet.png}
\vspace{-0.5em}
\caption{Mean per-decision wall-time across fleet sizes (log-scale) on
real Memphis OSM, three intensity panels. Each line is one HLP/LLP
combination averaged over 3 seeds. The hierarchical+learned-LLP filter
arm holds quality at the heuristic decomposition's level while compounding
the per-decision speedup.}
\label{fig:wall}
\end{figure}

\vspace{-0.5em}

\section{Composition with Uncertainty Modeling (Thrust~2 Barrier~1)}
The result above uses a deterministic travel-time oracle to establish the
dispatch-side speedup ceiling. The dispatcher abstracts travel times behind
a \texttt{TravelTimeOracle} interface, so the GNN-based stochastic
flow models proposed under Barrier~1 plug in as a drop-in replacement.
Two consequences follow. First, when uncertainty is introduced, the
LLP feasibility head must be re-trained on a stochastic-oracle dataset to
remain calibrated; the data-collection observer (which logs
$\sim$390k labeled candidates per sweep) is the same. Second, the
HLP signal that was insufficient in the deterministic regime grows
under stochasticity: rebalancing under uncertain travel times is
fundamentally a risk-management decision (allocate slack to regions whose
demand has high variance), which the heuristic does not do. We therefore
expect Barrier~1's GNN traffic model and Barrier~2's RL-based HLP to be
\emph{co-developed}: the stochastic oracle creates the headroom that
makes a learned HLP valuable.

\section{Position in the Multi-Modal Picture}
Within the broader Thrust~1 + Thrust~2 architecture, this work
operates the \textbf{microtransit leg} of the multi-modal mix: outskirt
$\to$ hub trips on capacitated on-demand vehicles. The same HLP/LLP
abstraction extends naturally to multi-modal dispatch by widening the
action space at both levels: the HLP allocates capacity across modes
(microtransit, fixed-route, vanpool, volunteer) given mode-specific cost
and reliability, while the LLP performs per-mode routing or transfer
scheduling. The decomposition argument we have validated in the
single-mode setting (per-decision tractability without quality loss)
is the precondition for that extension; multi-modal scheduling is
otherwise combinatorially intractable in real time.

\section{Digital-Twin Evaluation Plan (Thrust~2 Barrier~3)}
The digital twin (Co-PI Talusan, M2 milestone) is the validation
environment for these dispatchers under realistic operating conditions.
Concretely, we propose three evaluation phases against MATA / Feeonix /
Ford operational data:
\textbf{(P1) Calibration.} Replay one week of Feeonix microtransit traces
through the digital twin; tune the demand oracle and the OSM travel-time
distribution until the heuristic dispatcher reproduces observed assignment
rates and waiting times within $\pm$5\%.
\textbf{(P2) Stress testing.} Sweep fleet$\times$intensity$\times$incident
scenarios (traffic disruption, volunteer no-show, shift surge) and report
the speedup--quality Pareto frontier across HLP/LLP combinations. Success
metrics align with the proposal's \textit{Metrics of Success}: on-time
arrival reliability ($\geq$95\% on the heaviest cell), p95 per-decision
wall-time ($<$50\,ms at fleet=$10^3$), and passenger waiting time tails.
\textbf{(P3) Counterfactual evaluation.} For each historical incident,
compare the learned-LLP filter dispatcher's response against the recorded
operator response, quantifying both objective improvement and the
override rate by partner dispatchers (a proxy for Thrust~3 alignment).

\section{Status and Roadmap}
\textbf{Delivered:} HLP/LLP scaffold; OSM Memphis substrate; heuristic
HLP+LLP and learned-LLP plug-in (filter / stacked); 3-seed
$\times$ 9-cell sweeps; Hydra-driven master sweep over the four planner
combinations; HLP rebalance-decision data collector. \textbf{Next
(supporting Thrust~2 Year-1 milestones):} (a) wire the GNN stochastic
travel-time oracle (Barrier~1) into the dispatcher; (b) port a
multi-modal action space to the LLP insertion kernel; (c) reactivate
the learned-HLP path under stochastic travel and report the M2 digital-twin
calibration delta. The infrastructure built here directly supports the M2
(digital-twin validation) and M4 (operations prototype) milestones.

\end{document}
